% =============================================================================
% Section 4: Immutable Parent Pointers
% For: Recursive Spawning — BX3-Compliant Multi-Agent Delegation
% =============================================================================

\section{Immutable Parent Pointers}
\label{sec:immutable_parent_pointers}

The parent pointer chain is the central accountability primitive of the BX3 drift-prevention architecture. Unlike metadata that can be annotated, revised, or stripped after agent creation, a parent pointer in BX3 is cryptographically bound to the agent's identity record at the moment of provisioning and remains structurally inseparable from that identity for the agent's entire operational lifetime. This section formalizes the \ParentPointer{p}{c} relation and its immutability assignment invariant, the recursive chain construction algebra, the null-parent base case established by the root human, the Fact Layer's protocol-level enforcement of immutability, the Purpose Layer's spawn authorization gate, and the chain traversal algorithms that together constitute the RSP accountability substrate.

% ---------------------------------------------------------------------------- %
\subsection{ParentPointer: Formal Definition and Structural Properties}
\label{sec:pp_formal_definition}

\begin{definition}[ParentPointer]
\label{def:parent_pointer}
Let $\mathcal{A}$ be the set of all agents in a BX3-compliant system and let $\mathcal{H} \subset \mathcal{A}$ be the set of human principals. The \textbf{ParentPointer} relation $\ParentPointer{p}{c}$ holds \emph{iff} $p \in \mathcal{A}$ is the sole, immutable parent of $c \in \mathcal{A}$ at the moment of $c$'s provisioning. The functional notation $\parent(c) = p$ denotes the same relation. The inverse relation $\parent^{-1}(p) = \{\, c \in \mathcal{A} \mid \ParentPointer{p}{c} \,\}$ denotes the set of all direct children of $p$.
\end{definition}

The \ParentPointer{p}{c} relation satisfies four non-negotiable structural properties that collectively define the RSP parent-child relationship:

\begin{enumerate}
    \item[(P1) \textbf{Uniqueness}] For every child $c \in \mathcal{A}$, there exists exactly one $p \in \mathcal{A}$ such that $\ParentPointer{p}{c}$ holds at all times after provisioning. No child may have multiple simultaneous parent pointers. No mechanism exists by which an established pointer may be overwritten, redirected, or superseded by another. The Fact Layer's write protocol permits exactly one $\ParentPointer{p}{c}$ entry per child; a second entry for the same child is rejected as a protocol violation.

    \item[(P2) \textbf{Immutability after assignment}] After the provisioning write confirming $\ParentPointer{p}{c}$ is recorded in the Fact Layer ledger, no operation exists --- from any source, with any privilege level, under any system condition --- that can modify or delete $\ParentPointer{p}{c}$. The pointer is permanent for the operational lifetime of $c$. This is not a policy convention enforced by access control; it is a structural property of the write protocol: the modification operation is undefined in the Fact Layer's operation set. An agent $c$ cannot sever its parent relation post-creation, and no administrative principal can override this.

    \item[(P3) \textbf{Directionality}] The parent pointer is an intrinsic attribute of the child node, not a reference held by the parent. The child $c$ carries $\ParentPointer{p}{c}$ as an immutable attribute in its own Fact Layer state record. The parent $p$ holds no mutable reference to $c$. A parent that has terminated, crashed, or been decommissioned does not affect $c$'s ability to reference its parent. The chain is traversable from any node outward to its complete lineage independent of the operational state of any other node in the chain.

    \item[(P4) \textbf{Human Anchor Termination}] The root of any RSP chain is a human principal node $h \in \mathcal{H}$ for which $\ParentPointer{\bot}{h}$ holds, where $\bot$ denotes the null parent marker. This marker indicates that $h$ is the ultimate delegation origin. No machine actor may serve as a chain root; every chain terminates at a human principal by construction.
\end{enumerate}

\begin{invariant}[Immutability Assignment Invariant]
\label{inv:pp_immutability}
For any agent $c \in \mathcal{A}$, the parent relation $\parent(c)$ is assigned exactly once at the moment of $c$'s creation via a Fact Layer atomic write and sealed via SHA-256 hash chaining. This assignment is \textbf{permanent} and \textbf{immutable}. No subsequent operation may alter $\parent(c)$, nor may any actor --- including $c$ itself, the original parent $p$, any administrative principal, or the Purpose Layer post-authorization --- modify, clear, or dissociate this relation. The only terminal state for any agent $c$ is one in which $\parent(c)$ remains exactly as assigned at provisioning.
\end{invariant}

The immutability invariant has a direct consequence for agent identity itself. Because an agent's cryptographic identity record is bound to its sealed parent pointer via SHA-256 hash chaining, and because the parent pointer is immutable, the agent cannot exist in a state where its parent relation is undefined, modified, or replaced without this being immediately detectable as a chain integrity violation. The parent relation is not a property the agent \emph{has}; it is a constituent of what the agent \emph{is}. An agent attempting to disown its parent would need to break the SHA-256 hash chain binding its identity to that parent --- a computationally infeasible operation under standard cryptographic assumptions.

The null-parent case is reserved exclusively for the root human. For any machine agent $a \in \mathcal{A} \setminus \mathcal{H}$, $\parent(a) \neq \bot$ by construction. Any spawn event producing a machine agent with a null parent is rejected by the Fact Layer pre-commit hook. This enforces the human-grounding invariant: every machine agent in a BX3-compliant system has a traceable lineage to a human principal.

% ---------------------------------------------------------------------------- %
\subsection{Chain Construction Algebra}
\label{sec:chain_construction_algebra}

The parent pointer relation extends naturally to a complete ancestry chain through a recursive definition. The chain operator $\Chain{\cdot}$ is the primary tool for chain traversal, forensic verification, and accountability attribution.

\begin{definition}[Ancestry Chain]
\label{def:ancestry_chain}
For any agent $a \in \mathcal{A}$, the \textbf{ancestry chain} of $a$, denoted $\Chain(a)$, is defined recursively as:
\begin{equation}
\Chain(a) \;=\;
\begin{cases}
\emptyset & \text{if } \parent(a) = \bot \quad \text{(\textbf{base case: root human})} \\[12pt]
\{\,\ParentPointer{\parent(a)}{a}\,\} \; \cup \; \Chain(\parent(a)) & \text{if } \parent(a) \neq \bot \quad \text{(\textbf{recursive case})}
\end{cases}
\label{eq:chain_recursive_def}
\end{equation}
where $\cup$ denotes set union.
\end{definition}

Expanding the recursion one step makes the structure explicit:
\begin{equation}
\Chain(a) = \{\,\ParentPointer{\parent(a)}{a}\,\} \; \cup \; \bigl( \{\,\ParentPointer{\parent^{(2)}(a)}{\parent(a)}\,\} \; \cup \; \Chain(\parent^{(2)}(a)) \bigr).
\label{eq:chain_one_step_expansion}
\end{equation}

In fully expanded form, $\Chain(a)$ reveals the complete lineage as a set of parent pointer edges from $a$ back to the human anchor:
\begin{equation}
\Chain(a) = \bigl\{\, \ParentPointer{\parent(a)}{a},\; \ParentPointer{\parent^{(2)}(a)}{\parent(a)},\; \ParentPointer{\parent^{(3)}(a)}{\parent^{(2)}(a)},\; \ldots,\; \ParentPointer{\parent^{(n)}(a)}{\parent^{(n-1)}(a)} \,\bigr\}
\label{eq:chain_full_expansion}
\end{equation}
where $n$ is the smallest integer such that $\parent^{(n)}(a) = \bot$, and $\parent^{(n)}(a)$ is a human principal by (P4).

\begin{corol}[Chain Terminates at Human Anchor]
\label{corol:chain_human_termination}
For any agent $a \in \mathcal{A}$, $\Chain(a)$ always terminates at a human principal. Consequently, every agent in any valid RSP chain has a complete delegation lineage traceable to a human anchor, and no machine agent can serve as the root of a chain.
\end{corol}

\begin{examp}[Chain Construction]
Consider a delegation chain where human $H$ spawns agent $A$, $A$ spawns $B$, $B$ spawns $C$, and $C$ spawns $D$. The parent pointers established at each spawn event are:
\begin{align*}
\ParentPointer{H}{A} & \quad \text{(root $H$ spawns $A$)} \\
\ParentPointer{A}{B} & \quad \text{($A$ spawns $B$)} \\
\ParentPointer{B}{C} & \quad \text{($B$ spawns $C$)} \\
\ParentPointer{C}{D} & \quad \text{($C$ spawns $D$)}
\end{align*}
Applying the chain operator to $D$:
\begin{align*}
\Chain(D) &= \{\ParentPointer{C}{D}\} \; \cup \; \Chain(C) \\
          &= \{\ParentPointer{C}{D}\} \; \cup \; \{\ParentPointer{B}{C}\} \; \cup \; \Chain(B) \\
          &= \{\ParentPointer{C}{D},\; \ParentPointer{B}{C}\} \; \cup \; \{\ParentPointer{A}{B}\} \; \cup \; \Chain(A) \\
          &= \{\ParentPointer{C}{D},\; \ParentPointer{B}{C},\; \ParentPointer{A}{B}\} \; \cup \; \emptyset \\
          &= \{\ParentPointer{C}{D},\; \ParentPointer{B}{C},\; \ParentPointer{A}{B}\}.
\end{align*}
The chain $\Chain(D)$ contains exactly the three edges of the delegation chain from $D$ to the root human $H$. No edge is missing, no edge can be added after provisioning, and no edge can be modified without detection.
\end{examp}

The chain construction algebra satisfies four structural properties:

\begin{property}[Ancestry Chain Properties]
\label{prop:chain_algebra_properties}
For all agents $a, b \in \mathcal{A}$:
\begin{enumerate}
    \item \textbf{Well-definedness:} $\Chain(a)$ is always finite and always terminates at the root human. No infinite chains exist in a BX3-compliant system by construction, because the Bounds Engine enforces $D_{\max}$ at spawn time, providing a strict upper bound on chain depth.
    \item \textbf{Uniqueness:} If $b \in \Chain(a)$, then $\Chain(b) \subset \Chain(a)$. No agent appears twice in any chain.
    \item \textbf{Transitivity:} If $b \in \Chain(a)$ and $c \in \Chain(b)$, then $c \in \Chain(a)$. Chain membership is transitive.
    \item \textbf{No cycles:} There exists no agent $a$ such that $a \in \Chain(a)$. The \ParentPointer{} relation is a strict partial order (irreflexive, transitive, antisymmetric where defined).
\end{enumerate}
\end{property}

Property 4 (no cycles) is cryptographically enforced at spawn time: the Fact Layer pre-commit hook verifies that the proposed parent is not itself a descendant of the proposed child. Because $\parent$ is immutable once assigned, any cycle would require a back-reference to an ancestor that already has an immutable parent, which the pre-commit hook detects and rejects.

% ---------------------------------------------------------------------------- %
\subsection{The Root Human: Null Parent as Base Case}
\label{sec:root_human_base_case}

The root human occupies a structurally distinct position in the parent pointer chain. While all other agents carry a parent pointer referencing the agent that authorized their creation, the root human's $\parent$ relation is null by definition. This null-parent is not an exceptional condition requiring ad hoc handling; it is the recursive base case that terminates the chain construction algebra in Definition \ref{def:ancestry_chain}.

\begin{dfn}[Root Human]
\label{dfn:root_human_def}
The \textbf{root human} $h_r \in \mathcal{H}$ is the unique agent in the system such that $\parent(h_r) = \bot$. The root human is the \textbf{terminal accountability anchor}: all authority in the system flows from $h_r$, and all chains $\Chain(a)$ terminate at $h_r$. No agent may be created except by delegation from some existing agent; the chain must start at a human principal by architectural invariant.
\end{dfn}

The root human is the only agent for which the Fact Layer record contains a null parent marker $\bot$. This is a deliberate architectural decision: every chain must be grounded in human authorization. An agent whose chain terminates at anything other than a human principal is, by definition, unauthorized and must not enter operational state.

The Purpose Layer enforces this invariant at initialization and at every spawn event by verifying that the complete chain of any newly created agent terminates at a human principal in $\mathcal{H}$. The Fact Layer independently enforces it by rejecting any spawn event that would produce a machine agent with a null parent.

This design reflects a foundational accountability principle: \textbf{no agent may authorize the existence of another agent without a traceable chain of human delegation}. The root human is the embodiment of this principle at the base of every chain. It is the singular point of original authority from which all legitimate agency flows, and it is the terminal reference point for all forensic audit and drift detection procedures.

% ---------------------------------------------------------------------------- %
\subsection{Immutability Enforcement: The Fact Layer Write Protocol}
\label{sec:fact_layer_enforcement}

The immutability of the parent pointer is not a policy convention; it is a structural property of the Fact Layer's write protocol. The Fact Layer is the append-only forensic ledger that records all agent creation events, seals each record with a SHA-256 hash binding it to the prior record, and exposes no modification interface for any established parent pointer.

At the moment of agent creation, the Fact Layer records the following data into a cryptographically sealed block:
\begin{itemize}
    \item The identity hash of the newly created agent $c$
    \item The identity hash of the parent agent $p$
    \item The complete purpose declaration of $c$
    \item The authority bounds assigned to $c$
    \item The parent pointer relation $\ParentPointer{p}{c}$
    \item A timestamp and sequence number
    \item The hash of the immediately preceding Fact Layer block
\end{itemize}

This block is hashed with SHA-256 to produce a commitment cryptographically bound to the entire prior chain. Any attempt to modify the parent pointer after sealing --- whether by altering $p$, setting $\parent(c)$ to null, or substituting a different parent --- would require recomputing the hash of the block containing $c$, which would invalidate the hash of every subsequent block in the chain. The SHA-256 second-preimage resistance property ensures that there exists no computationally feasible method to produce a modified chain that hashes to the same root commitment without breaking the cryptographic primitive.

The Fact Layer enforces immutability through four distinct mechanisms that operate at the protocol level:

\begin{enumerate}
    \item[(W1) \textbf{No modification operation defined.}] The Fact Layer operation set contains no entry for $\mathrm{modify}(\ParentPointer{p}{c})$. Any request to alter, overwrite, redirect, or delete an established parent pointer is rejected as an unrecognized operation before it reaches the ledger. This is not a policy response to a recognized request; it is a protocol-level rejection of an entire operation class.

    \item[(W2) \textbf{Unconditional rejection.}] Any request to alter, overwrite, redirect, or delete $\ParentPointer{p}{c}$ is rejected before the Fact Layer attempts to execute it. The rejection is unconditional: it does not depend on the privilege level of the requesting actor, the system state, or any configurable policy. All sources receive identical treatment. There is no privileged actor capable of overriding this rejection.

    \item[(W3) \textbf{Sealed memory envelope.}] The parent pointer field $\ParentPointer{p}{c}$ is encrypted at rest using a Fact Layer-only symmetric key. Even a compromised system component cannot read $\ParentPointer{p}{c}$ because it never has access to the decryption key. The envelope carries an HMAC-SHA-256 signature computed over its contents; any modification to the ciphertext invalidates the signature. The Fact Layer decrypts the envelope only when reading $\ParentPointer{p}{c}$ to service a chain-traversal or audit request, and re-encrypts and re-signs after each authorized read.

    \item[(W4) \textbf{Atomic provisioning.}] The parent pointer $\ParentPointer{p}{c}$ and its Purpose Layer warrant $\phi$ are created as a single atomic Fact Layer write. There is no temporal window in which the pointer exists without authorization, or in which authorization exists without a corresponding pointer. This eliminates the class of exploits where a parent pointer is first established and then retroactively authorized by a subsequent warrant.
\end{enumerate}

The distinction between access-control-based immutability and protocol-level immutability is the distinction between a promise and a guarantee. In an access-control model, immutability is enforced by restricting which principals may issue modification operations. If an attacker acquires sufficient privileges --- through credential theft, insider compromise, or software vulnerability --- the access control barrier is circumvented and immutability is violated. The Fact Layer write protocol eliminates this entire failure class: retroactive manipulation of chain topology for the purpose of accountability laundering is structurally impossible, because the modification operation is not defined in the protocol.

The immutability guarantee exhibits four properties that distinguish it from any access-control-based approach:

\begin{enumerate}
    \item[(E1) \textbf{No override path.}] There is no configuration flag, emergency pathway, or administrative override capable of modifying $\ParentPointer{p}{c}$ after provisioning. The Fact Layer write protocol has no analog to a system administrator override or a break-glass mechanism.

    \item[(E2) \textbf{Failure-state resilience.}] The immutability guarantee holds even under system conditions that disable other enforcement mechanisms --- including catastrophic failure, kernel compromise, or memory corruption in other system components. The parent pointer is stored in non-volatile, sealed Fact Layer memory architecturally isolated from the general-purpose execution environment.

    \item[(E3) \textbf{Verification independence.}] Any agent or external auditor with read access to the Fact Layer can independently recompute the SHA-256 hash chain from the root human forward and compare the result against the stored root commitment. If any discrepancy is found, the chain is declared invalid and the agent is suspended. No trusted intermediary whose certification is required to verify chain integrity.

    \item[(E4) \textbf{Atomic warrant-pointer coupling.}] The parent pointer and its Purpose Layer warrant are created as a single atomic Fact Layer write. There is no temporal window in which $\ParentPointer{p}{c}$ is established without a matching warrant, or in which a warrant exists without a corresponding parent pointer. This coupling ensures that every sealed parent pointer is always accompanied by a valid Purpose Layer authorization at the moment of creation.
\end{enumerate}

% ---------------------------------------------------------------------------- %
\subsection{Spawn Authorization: The Purpose Layer Role}
\label{sec:purpose_layer_spawn_auth}

While the Fact Layer provides cryptographic immutability for the parent pointer, the Purpose Layer provides the authorization mechanism that determines which parent-child relationships are permitted to be created in the first place. Spawn authorization is the gate that ensures every $\ParentPointer{p}{c}$ relation in the system represents a legitimate, scope-refined delegation of authority from a human principal downward.

The spawn authorization procedure $P_{\mathrm{spawn}}$ evaluates three mandatory conditions for any proposed spawn event:

\begin{enumerate}
    \item[(S1) \textbf{Scope refinement.}] The child scope must be a strict subset of the parent scope: $R(c) \subsetneq R(p)$. No lateral expansion, no superset, no equality. This ensures each spawn event narrows the delegation scope, preserving the intent-refinement property of the chain. A child may not receive broader authority than its parent possesses.

    \item[(S2) \textbf{Operational necessity.}] The parent must declare, in the Fact Layer authorization ledger, the precise condition requiring child spawning. The declaration must identify why the condition cannot be resolved by the parent acting alone, and why the scope of the proposed child is the minimum scope sufficient to resolve it. This prevents spawn events that inflate the delegation chain without corresponding operational justification.

    \item[(S3) \textbf{Human intent alignment.}] The proposed child scope must be contained within the declared purpose of the parent, traceable through the chain to the root human's purpose clause. This ensures no spawn event introduces operational intent that was not present in the root human's original authorization. The child cannot inherit a purpose that its parent does not itself possess through the chain.
\end{enumerate}

The mandatory sequencing is:
\begin{enumerate}
    \item The parent $p$ submits a spawn request to the Purpose Layer specifying the proposed child scope $\sigma_c$ and the operational necessity declaration.
    \item The Purpose Layer evaluates (S1), (S2), and (S3). If any condition fails, the spawn request is rejected and the Fact Layer records a \texttt{spawn-refused}$(p, c, \text{reason})$ event.
    \item If all conditions hold, the Purpose Layer issues a cryptographic warrant $\phi$ certifying that the spawn is authorized, and transmits the warrant to the Fact Layer.
    \item The Fact Layer receives the warrant and records $\ParentPointer{p}{c}$ as an immutable attribute of $c$.
\end{enumerate}

The Fact Layer will not complete the provisioning write for a spawn event without a valid Purpose Layer warrant. This coupling ensures that every parent pointer in the system is both authorized (Purpose Layer) and sealed (Fact Layer). No child node $c$ can have an established $\ParentPointer{p}{c}$ in the Fact Layer ledger unless the Purpose Layer has issued a valid warrant $\phi_c$ for the spawn event.

The Purpose Layer's spawn authorization is what gives the parent pointer chain its semantic meaning beyond mere structural data. The parent pointer is not simply a reference from child to parent; it is a cryptographically sealed record that the child was authorized to exist by the parent, that the parent's authority was sufficient to delegate the requested purpose, and that the entire chain is grounded in human authorization. Every $\ParentPointer{p}{c}$ edge is therefore not only immutable but also always legitimate at the moment of its creation.

% ---------------------------------------------------------------------------- %
\subsection{Chain Traversal Algorithm}
\label{sec:chain_traversal_algo}

The ancestry chain is a queryable data structure enabling any authorized actor to reconstruct the complete delegation path of an agent at any point during execution. The Chain Query Interface (CQI) exposes this capability as a first-class forensic operation.

\begin{algorithm}[htbp]
\caption{Chain-Traverse($n$)}
\label{alg:chain_traverse}
\begin{algorithmic}[1]
\Require Node $n \in \mathcal{A}$
\Ensure Returns $\Chain{n}$ as the ordered set of all $\ParentPointer{}$ edges from $n$ to the root human
\State $\text{current} \gets n$
\State $\text{chain} \gets \emptyset$
\Loop
    \State $p \gets \parent(\text{current})$
    \If{$p = \bot$}
        \State \Return $\text{chain}$ \Comment*{Base case: root human reached}
    \EndIf
    \State $\text{chain} \gets \text{chain} \cup \{\,\ParentPointer{p}{\text{current}}\,\}$
    \State $\text{current} \gets p$
\EndLoop
\end{algorithmic}
\end{algorithm}

Chain-Traverse follows parent pointers iteratively from the target node until reaching the root human, collecting each edge in the chain set. The algorithm terminates because the chain is guaranteed to be finite: the Bounds Engine enforces a maximum delegation depth $D_{\max}$ at spawn time, providing a strict upper bound on chain length, and the no-cycles invariant (Property \ref{prop:chain_algebra_properties}) ensures that the traversal cannot loop indefinitely.

\begin{algorithm}[htbp]
\caption{Chain-Verify($n$)}
\label{alg:chain_verify}
\begin{algorithmic}[1]
\Require Node $n \in \mathcal{A}$
\Ensure Returns \texttt{true} if $\Chain{n}$ satisfies all RSP structural constraints; $\texttt{false}(i)$ with index $i$ of first violation otherwise
\State $chain \gets \text{Chain-Traverse}(n)$
\State $m \gets |chain|$
\Statex
\Comment*{V1: Verify human anchor termination}
\State $root \gets chain[m-1]$
\If{$\neg \isHuman(root)$}
    \State \Return $\texttt{false}$ \Comment*{Chain does not terminate at human anchor}
\EndIf
\Statex
\Comment*{V2: Verify Purpose Layer warrant at each link}
\For{$i \leftarrow 0$ \To $m-1$}
    \State $child \leftarrow chain[i].child$
    \State $parent \leftarrow chain[i].parent$
    \If{$\neg \warrant\_exists(child, parent)$}
        \State \Return $\texttt{false}$ \Comment*{Missing Purpose Layer warrant for spawn}
    \EndIf
\EndFor
\Statex
\Comment*{V3: Verify scope refinement at each link}
\For{$i \leftarrow 0$ \To $m-2$}
    \State $child \leftarrow chain[i].child$
    \State $parent \leftarrow chain[i+1].parent$
    \If{$\neg (R(child) \subsetneq R(parent))$}
        \State \Return $\texttt{false}$ \Comment*{Scope refinement constraint violated}
    \EndIf
\EndFor
\Statex
\Comment*{V4: Verify depth bound compliance}
\For{$i \leftarrow 0$ \To $m-1$}
    \State $\text{node} \gets chain[i].child$
    \If{$\mathrm{depth}(\text{node}) > D_{\max}$}
        \State \Return $\texttt{false}$ \Comment*{Depth bound exceeded}
    \EndIf
\EndFor
\State \Return $\texttt{true}$
\end{algorithmic}
\end{algorithm}

Chain-Verify returns \texttt{true} if and only if all four verification conditions hold simultaneously. V1 confirms the chain terminates at a human anchor. V2 confirms every spawn was authorized by the Purpose Layer with a valid warrant. V3 confirms scope refinement was respected at every link. V4 confirms the depth bound was respected throughout. A chain that passes Chain-Verify is a fully authorized delegation chain terminating at a human principal.

\begin{prop}[Chain-Verify Soundness]
\label{prop:chain_verify_soundness}
If Chain-Verify($n$) returns \texttt{true}, then $\Chain{n}$ is a valid RSP delegation chain terminating at a human principal, and all spawn events in the chain satisfy the Purpose Layer authorization conditions (S1), (S2), and (S3).
\end{prop}

\begin{prop}[Chain-Verify Completeness]
\label{prop:chain_verify_completeness}
If $\Chain{n}$ is a valid RSP delegation chain terminating at a human principal, then Chain-Verify($n$) returns \texttt{true}.
\end{prop}

\textbf{Complexity.} Chain-Traverse runs in $O(d)$ time where $d$ is the chain depth, bounded by $d \leq D_{\max} + 1$. Chain-Verify runs in $O(m)$ for the traversal plus $O(m)$ for each of three verification loops, yielding $O(m)$ total where $m$ is the chain length, also bounded by $D_{\max} + 1$. Both algorithms are linear in chain length and independent of the total number of nodes in the system. Mutable parent pointers would allow the chain to be extended retroactively, making the effective depth potentially unbounded and the worst-case traversal cost unbounded. Immutability is precisely what makes $D_{\max}$ a structural upper bound on traversal cost.

The CQI exposes three primary query operations to all authorized agents and human overseers:
\begin{itemize}
    \item $\proc{GetAncestors}(a)$: Returns the ordered list $[a_0, a_1, \ldots, a_n]$ where $a_n$ is the root human, via Chain-Traverse.
    \item $\proc{GetDepth}(a)$: Returns the integer depth of agent $a$ --- the number of delegation hops from the root human to $a$ --- computed as $|\Chain(a)|$.
    \item $\proc{VerifyChain}(a)$: Wrapper around Chain-Verify that returns $\mathrm{VALID}$ if integrity holds or $\mathrm{INVALID}(i)$ with the index of the first invalid block if it does not.
\end{itemize}

The $\proc{VerifyChain}$ operation is the forensic workhorse: it enables any party with read access to the Fact Layer to independently verify that the ancestry chain of any agent has not been tampered with, without requiring trust in the agent itself or any intermediate party.

% ---------------------------------------------------------------------------- %
\subsection{Connection to Drift Prevention}
\label{sec:drift_prevention_link}

The immutable parent pointer chain is the structural foundation on which the drift prevention guarantee is built. Its role operates at three distinct levels:

\begin{enumerate}
    \item \textbf{Structural Immutability.} Because the parent pointer is sealed at creation and cannot be modified, the chain is permanently legible. Any attempt to modify a chain post-creation is detectable by the Fact Layer's hash chain verification. An agent cannot dissociate itself from its ancestry to escape accountability.

    \item \textbf{Intent Preservation.} The parent pointer chain preserves the full delegation path, enabling reconstruction of the original purpose at any depth. Even if a sub-agent at depth $n$ exhibits behavior that appears drifted, the chain enables an auditor to trace the purpose sub-objectives from the root human forward and identify where intent propagation diverged from the authorized path.

    \item \textbf{Forensic Completeness.} The combination of sealed parent pointers, hash chain integrity, and the CQI means that the complete execution provenance of any agent is always available and always verifiable. This is the prerequisite for forensic drift detection procedures that complement the structural drift prevention provided by the Purpose Layer and the runtime enforcement provided by the Truth Gate.
\end{enumerate}

These three levels map directly onto the three layers of the BX3 architecture: the Fact Layer provides structural immutability and forensic completeness; the Purpose Layer provides intent preservation through spawn authorization and the human ground invariant; and the Truth Gate provides runtime enforcement confirming that any action taken by any agent is consistent with its authorized purpose and its sealed chain. The parent pointer chain is thus not merely a data structure --- it is the connective tissue linking the three layers into a coherent drift-prevention system. Within this architecture, the drift triad (orphaned agents, drifted agents, and simultaneously operating unauthorized agents) is not merely statistically unlikely; it is structurally excluded from the set of architecturally possible states.